\documentclass[11pt,a4paper]{article}
\usepackage[utf8]{inputenc}
\usepackage[T1]{fontenc}
\usepackage[italian]{babel}
\usepackage{geometry}
\usepackage{amsmath}
\usepackage{booktabs}
\geometry{margin=1in}
\setlength{\parskip}{0.4em}

\title{Note di lavoro: cosa rende un grafo ``difficile''\\ per imparare la connettività}
\author{}
\date{\today}

\begin{document}
\maketitle

Questo documento spiega, passo per passo e in parole semplici, il ragionamento
dietro i prossimi esperimenti. Niente gergo non spiegato.

\section{Due domande diverse che stiamo inseguendo}

Stiamo facendo due cose distinte. È importante non confonderle.

\paragraph{Domanda A --- ``possiamo fare un modello migliore?''}
Lo standard transformer (quello del paper) generalizza male. La domanda A è:
\emph{se cambiamo un pezzo dell'architettura, otteniamo un modello che sbaglia
meno?} Il pezzo che stiamo cambiando è il \textbf{read-out} (lo spiego tra poco).
Questa domanda è ``costruttiva'': vogliamo un modello migliore.

\paragraph{Domanda B --- ``cosa rende un grafo difficile?''}
Il paper sostiene che la difficoltà di un grafo sia il suo \textbf{diametro} (la
distanza massima tra due nodi). La domanda B è: \emph{è davvero il diametro? O
c'è una proprietà del grafo che predice meglio quando il modello sbaglia?}
Questa domanda è ``descrittiva'': vogliamo capire, non costruire.

\paragraph{Come si parlano.} Sono separate, ma collegate: se con la domanda B
capiamo \emph{perché} certi grafi sono difficili, allora sappiamo \emph{quale}
modifica provare nella domanda A. Esempio: se scopriamo che la difficoltà sono i
``colli di bottiglia'' del grafo, allora un'architettura che aiuta a superarli è
la cosa giusta da costruire.

\section{Come funziona il modello, in parole semplici}

\begin{itemize}
  \item \textbf{Input}: la matrice di adiacenza (chi è collegato a chi).
  \item Il transformer costruisce, per ogni nodo, un \textbf{embedding}: un
        vettore di numeri che ``riassume'' quel nodo.
  \item Il \textbf{read-out} è la regola che, guardando gli embedding di due
        nodi $h_i$ e $h_j$, decide: ``sono nella stessa componente, sì o no?''
\end{itemize}

Usiamo due read-out diversi:
\begin{itemize}
  \item \textbf{Lineare} (quello standard): una formula generica che combina i
        due vettori. Non richiede che i due vettori si somiglino.
  \item \textbf{Similarità}: dice ``connessi'' se i due vettori
        \textbf{puntano nella stessa direzione} (coseno alto tra $h_i$ e $h_j$).
\end{itemize}

\section{L'idea chiave: con la similarità, ``connesso = stesso vettore''}

Se uso il read-out a \textbf{similarità}, il modello ideale deve fare questo:
dare a tutti i nodi della \emph{stessa} componente un embedding (quasi)
identico, e a nodi di componenti \emph{diverse} embedding diversi.

\paragraph{Analogia.} Pensa a colorare i nodi. Ogni componente connessa è un
colore. Il modello ideale dà lo stesso colore a tutti i nodi collegati tra loro.
Allora ``$i$ e $j$ connessi?'' diventa semplicemente ``hanno lo stesso colore?''.

Attenzione: questo vale \emph{solo} con il read-out a similarità. Con il read-out
lineare il modello \emph{non} ha bisogno di rendere uguali gli embedding, quindi
guardare ``quanto sono simili i vettori'' lì non ci dice nulla di utile. Questo
è un punto importante e ci torneremo.

\section{Perché alcuni grafi sono difficili: l'informazione deve ``viaggiare''}

Come fa il modello a dare lo stesso embedding (colore) a due nodi connessi ma
\emph{lontani}? L'informazione deve propagarsi lungo il grafo, di vicino in
vicino. \textbf{Ogni layer del transformer è un passo di propagazione}: come una
voce che si passa di persona in persona.

Il nostro modello ha 2 layer, quindi fa circa 2 passi. Allora:
\begin{itemize}
  \item Se due nodi connessi sono ``facili da raggiungere'' (tanti cammini, ben
        collegati), 2 passi bastano a dargli lo stesso colore $\rightarrow$ facile.
  \item Se sono separati da un \textbf{collo di bottiglia} (un solo ponte
        stretto) o da un lungo corridoio, 2 passi \emph{non} bastano $\rightarrow$
        il modello non riesce a dargli lo stesso colore $\rightarrow$ sbaglia.
\end{itemize}

Nei modelli su grafi questo fenomeno ha un nome: \textbf{over-squashing}
(l'informazione si ``strozza'' nei colli di bottiglia).

\paragraph{Esempi concreti (dai nostri dati reali).}
\begin{center}
\begin{tabular}{lccc}
\toprule
grafo & com'è fatto & diametro & quanto sbaglia (exact) \\
\midrule
expander  & tutti ben collegati          & 5  & 0.99 (facile) \\
clique    & tutti con tutti              & 1  & facile (sul reach) \\
catena lunga (1chain) & fila indiana     & 19 & 0.45 \\
barbell   & due gruppi $+$ ponte stretto & 11 & 0.22 (difficile) \\
\bottomrule
\end{tabular}
\end{center}

Guarda l'ultimo confronto: il \textbf{barbell} ha diametro \emph{più piccolo}
(11) della catena lunga (19), eppure è \emph{più difficile} (0.22 contro 0.45)!
Quindi il diametro \emph{non} spiega la difficoltà: il vero problema del barbell
è il ponte stretto, non la distanza. Questo è il cuore della domanda B.

\section{Come misurare la difficoltà: tre candidati}

\paragraph{(a) Diametro --- la proposta del paper.}
La distanza massima (in numero di passi) tra due nodi. \emph{Limite}: conta solo
i passi, non \emph{quanti} cammini ci sono. Due nodi a 5 passi lungo un corridoio
(un solo cammino) sono molto più difficili di due nodi a 5 passi dentro un gruppo
denso (mille cammini). Il diametro non vede questa differenza.

\paragraph{(b) Spectral gap --- ``quanto in fretta si mescola il grafo''.}
Un singolo numero per tutto il grafo. Grande = l'informazione si mescola in fretta
(facile); piccolo = c'è un collo di bottiglia (difficile). È utile, ma è un solo
numero per l'intero grafo, non per la singola coppia di nodi.

\paragraph{(c) Resistenza effettiva --- per ogni coppia di nodi.}
Immagina il grafo come una rete di fili elettrici, ogni arco è una resistenza.
La ``resistenza effettiva'' tra due nodi misura quanto è difficile far passare
corrente tra loro.
\begin{itemize}
  \item Tanti cammini paralleli $\rightarrow$ bassa resistenza $\rightarrow$ facile.
  \item Un solo cammino lungo, o un collo di bottiglia $\rightarrow$ alta
        resistenza $\rightarrow$ difficile.
\end{itemize}
\emph{Analogia}: due città collegate da 10 autostrade contro due città collegate
da 1 strada sterrata. La ``distanza'' può essere la stessa, ma la difficoltà di
spostarsi è diversa. La resistenza coglie questa differenza, il diametro no. Ed è
una misura \textbf{per coppia di nodi}, quindi più fine del diametro.

\paragraph{(d) $P^L$ --- la più fedele al nostro modello.}
$P$ vuol dire ``fai un passo casuale verso un vicino''. $P^L$ vuol dire ``dopo
$L$ passi, quanta probabilità di essere arrivato dall'altro nodo''. Mettiamo
$L = 2$ perché il modello ha 2 layer (= 2 passi). Se dopo 2 passi non ci si
arriva quasi mai, allora un modello a 2 layer \emph{non può} aver collegato quei
due nodi $\rightarrow$ coppia difficile. Il bello di questa misura è che usa
esattamente il ``budget di passi'' del modello (cioè il numero di layer).

\section{L'esperimento che voglio fare}

Prendo i modelli già allenati. Per ogni coppia di nodi \emph{connessi}, raccolgo:
\begin{enumerate}
  \item quanto il modello ha sbagliato su quella coppia;
  \item (solo per il read-out a similarità) quanto sono diventati simili i loro
        embedding --- cioè se il modello è riuscito a dargli lo stesso ``colore'';
  \item le tre misure di difficoltà: diametro, resistenza effettiva, $P^L$.
\end{enumerate}

Poi guardo: \textbf{quale misura predice meglio l'errore?} Se la resistenza (o
$P^L$) batte il diametro, abbiamo trovato la nozione giusta di difficoltà --- e
questo è esattamente il punto delle domande B (capire) e A (sapere cosa migliorare).

\section{Le due scelte che ti chiedevo (spiegate bene)}

\subsection*{Scelta 1: pesi ``uniformi'' o pesi ``veri'' del modello?}

Quando il modello propaga l'informazione, \emph{non} dà lo stesso peso a tutti i
vicini: il meccanismo di attention decide a chi ``dare retta'' di più, e questi
pesi li impara durante il training. Quando calcolo $P^L$ (la propagazione) ho due
strade:
\begin{itemize}
  \item \textbf{Uniforme} (semplice): faccio finta che ogni nodo tratti tutti i
        vicini allo stesso modo. Dipende solo dal grafo, non dal modello. Facile
        e veloce da calcolare.
  \item \textbf{Pesi veri}: uso i veri pesi di attention che il modello ha
        imparato. Più fedele a ciò che il modello fa davvero, ma più lavoro.
\end{itemize}

La frase ``proxy falsificabile'' che non ti diceva niente significa questo:
\emph{proxy} = approssimazione semplice (la versione uniforme); \emph{falsificabile}
= fa una previsione netta (``le coppie con propagazione bassa sono quelle che il
modello sbaglia'') che possiamo verificare e, se è sbagliata, scoprirlo subito.

\textbf{Proposta}: parto dalla versione \emph{uniforme} perché è semplice. Se
funziona, ci fermiamo lì. Se non basta a spiegare gli errori, allora passo ai
\emph{pesi veri} del modello. In breve: non complico finché non serve.

\subsection*{Scelta 2: come confronto la resistenza tra grafi diversi?}

La resistenza ha valori su scale diverse a seconda di quanto è grande il grafo,
quindi confrontare i numeri ``grezzi'' tra grafi diversi è ingannevole. Soluzione:
invece dei valori assoluti guardo l'\textbf{ordine} (quale coppia è più difficile
di quale) e confronto gli ordini. Così la scala non mi inganna.

\section{Un'ultima distinzione: ``reach'' contro ``cut''}

Ci sono due compiti diversi dentro la connettività:
\begin{itemize}
  \item \textbf{Reach}: rendere \emph{uguali} due nodi connessi ma lontani. La
        difficoltà qui è la resistenza / il collo di bottiglia (tutto quanto detto
        sopra).
  \item \textbf{Cut}: rendere \emph{diversi} due nodi di componenti separate.
        Strutturalmente è facile (sono scollegati). Quando il modello sbaglia qui
        (es. sulle due cricche), è per un altro motivo --- una scorciatoia tipo
        ``ha tanti vicini, quindi lo dico connesso'' --- non per un collo di
        bottiglia.
\end{itemize}
Per questo, nell'analisi, tengo separate le coppie ``reach'' dalle coppie ``cut'':
le misure di difficoltà di cui sopra spiegano il \emph{reach}; il \emph{cut} è un
problema diverso che va guardato a parte.

\section*{In sintesi}

\begin{enumerate}
  \item Con il read-out a similarità, il modello ideale dà lo stesso ``colore''
        (embedding) a nodi connessi. Sì, la tua intuizione era giusta.
  \item Alcuni grafi sono difficili perché l'informazione fatica a viaggiare tra
        nodi connessi (colli di bottiglia) --- è l'over-squashing.
  \item Il diametro non basta a misurarlo (il barbell lo dimostra). Candidati
        migliori: spectral gap, resistenza effettiva, $P^L$.
  \item L'esperimento: vedere quale misura predice gli errori del modello.
  \item Le due scelte: parto dalla versione semplice (uniforme) e uso l'ordine
        invece dei valori grezzi.
\end{enumerate}

\end{document}
